\subsection{Les G-cartes}

\begin{figure*}[h!]
    \begin{subfigure}{.19\textwidth}
        \centering
        \includegraphics[height=.18\textheight]{decomposition-3D-no-decomp}
        \caption{}
        \label{fig:demo-fig-ABCDE}
    \end{subfigure}
    \hfill
    \begin{subfigure}{.19\textwidth}
        \centering
        \includegraphics[height=.18\textheight]{decomposition-3D-3-cell}
        \caption{}
        \label{fig:decomp-volume}
    \end{subfigure}
    \hfill
    \begin{subfigure}{.19\textwidth}
        \centering
        \includegraphics[height=.18\textheight]{decomposition-3D-2-cell}
        \caption{}
        \label{fig:demo-fig-faces-decomposition}
        \end{subfigure}
    \hfill
    \begin{subfigure}{.19\textwidth}
        \centering
        \includegraphics[height=.18\textheight]{decomposition-3D-1-cell}
        \caption{}
        \label{fig:demo-fig-show-1-links}
    \end{subfigure}
    \hfill
    \begin{subfigure}{.19\textwidth}
        \centering
        \includegraphics[height=.18\textheight]{decomposition-3D-0-cell}
        \caption{}
        \label{fig:demo-fig-show-all-links}
    \end{subfigure}
    \caption{Décomposition d'un modèle géométrique 3D par ses cellule}
    \label{fig:demo-fig-cell-decomposition}
\end{figure*}

La G-carte est une classe particulière de graphe étiqueté permettant de
représenter un modèle de dimension $n$ à partir de sa structure topologique,
c'est-à-dire par ses volumes, faces, arêtes, sommets et les relations entre
ceux-ci pour une G-carte de dimension 3. %
Prenons l'exemple d'un objet géométrique 3D (Figure~\ref{fig:demo-fig-ABCDE})
composé de deux volumes (Figure~\ref{fig:decomp-volume}) reliés entre eux par
des $3$-liaisons (vert). %
Chaque volume est composé de faces
(Figure~\ref{fig:demo-fig-faces-decomposition}) reliées entre elles par des
$2$-liaisons (bleu). %
Ces faces sont composées de sommets (Figure~\ref{fig:demo-fig-show-1-links})
reliés entre eux par des $1$-liaisons (rouge). %
Et ces sommets sont eux-mêmes reliés entre par des brins
(Figure~\ref{fig:demo-fig-show-all-links}) reliés par des $0$-liaisons (noir). %

\begin{figure*}[h!]
  \begin{subfigure}{.19\textwidth}
    \centering
    \includegraphics[height=.18\textheight]{orbit-3D-vertex-123.pdf}
    \caption{}
    \label{fig:orbit-3D-vertex-123}
  \end{subfigure}
  \hfill
  \begin{subfigure}{.19\textwidth}
    \centering
    \includegraphics[height=.18\textheight]{orbit-3D-edge-023.pdf}
    \caption{}
    \label{fig:orbit-3D-edge-023}
  \end{subfigure}
  \hfill
  \begin{subfigure}{.19\textwidth}
    \centering
    \includegraphics[height=.18\textheight]{orbit-3D-face-013}
    \caption{}
    \label{fig:orbit-3D-face-013}
  \end{subfigure}
  \hfill
  \begin{subfigure}{.19\textwidth}
    \centering
    \includegraphics[height=.18\textheight]{orbit-3D-volume-012}
    \caption{}
    \label{fig:orbit-3D-volume-012}
  \end{subfigure}
  \hfill
  \begin{subfigure}{.19\textwidth}
    \centering
    \includegraphics[height=.18\textheight]{orbit-3D-edge-0}
    \caption{}
    \label{fig:orbit-3D-edge-02}
  \end{subfigure}
  \caption{Décomposition d'un modèle géométrique 3D par ses orbites}
  \label{fig:orbit-decomposition}
\end{figure*}

Dans une G-carte, un brin est l'entité de plus petite dimension. %
Ceux-ci sont donc utilisés au sein de ce que l'on appelle des orbites afin de
représenter des entités topologiques. %
Une orbite est composée d'un ensemble des brins et d'un ensemble de liaisons
dont les étiquettes décrivent un type d'orbite $\langle o \rangle$. %
Ainsi, une orbite est l'ensemble des brins atteignables à partir d'un brin $b$
et des $i$-liaisons parcourues tel que $i$ appartient à $\langle o \rangle$. %
Prenons l'exemple des orbites illustrées dans la
Figure~\ref{fig:orbit-decomposition} et incidentes au brin $a$. %
Le sommet (Figure~\ref{fig:orbit-3D-vertex-123}), de type
$\langle 1,2,3 \rangle$, est composé des brins atteignables par les liaisons
étiquetées $1$, $2$ et $3$. %
L'arête (Figure~\ref{fig:orbit-3D-edge-023}), de type $\langle 0,2,3 \rangle$
est composée des brins atteignables par les liaisons étiquetées $0$, $2$ et
$3$. %
La face (Figure~\ref{fig:orbit-3D-face-013}), est de type
$\langle 0,1,3 \rangle$ et le volume (Figure~\ref{fig:orbit-3D-volume-012}) de
type $\langle 0,1,2 \rangle$. %
Les orbites permettent également de représenter des entités topologiques de
dimensions inférieures, par exemple l'arête de volume
(Figure~\ref{fig:orbit-3D-edge-02}) de type $\langle 0,2 \rangle$. %

\subsection{Les règles de transformation de graphe}

Les règles de transformations de graphe Jerboa permettent la formalisation des
opérations sur des G-cartes. %

En quelques mots, une règle $r: L \longrightarrow R$ et un filtrage $m: L \to G$
vers une G-carte $G$, décrivent la transformation $t: G \longrightarrow^{r,m} H$
de $G$ vers $H$. %
Le filtrage $m$ permet de remplacer un sous-graphe de $G$, décrit par le membre
gauche $L$ de la règle, par un autre sous-graphe décrit par le membre droit $R$
et produire le graphe $H$. %

\begin{figure}[h!]
  \centering
  \begin{subfigure}[t]{.32\linewidth}
    \includegraphics[width=.8\textwidth]{demo-fig-extrusion-left-match}
    \caption{G}
    \label{fig:face-to-extrude}
  \end{subfigure}
  \hfill
  \begin{subfigure}[t]{.32\linewidth}
    \includegraphics[width=.8\textwidth]{demo-fig-extrusion-right-match2}
    \caption{H}
    \label{fig:extruded-cube}
  \end{subfigure}
  \hfill
  \begin{subfigure}[t]{.32\textwidth}
  \centering
    \includegraphics[width=.8\textwidth]{triangulated-cube-article-view.pdf}
    \caption{I}
    \label{fig:demo-cube-triangulated}
  \end{subfigure}
  \caption{Extrusion d'une face}
  \label{fig:demo-sequence}
\end{figure}

Présentons ces règles avec l'exemple de la séquence suivante : nous créons une
face (Figure~\ref{fig:face-to-extrude}), cette face est extrudée pour produire
un cube (Figure~\ref{fig:extruded-cube}) et enfin, nous triangulons la face avant gauche de ce cube (Figure~\ref{fig:demo-cube-triangulated}). %

\begin{figure}[h!]
  \centering
  \includegraphics[height=.2\textheight]{rule-create-square-face}
  \caption{Règle de création de face}
  \label{fig:rule-create-square-face}
\end{figure}

La règle qui crée une face carrée (Figure~\ref{fig:rule-create-square-face}) a
son membre gauche vide et son membre droit composé des nœuds $n_{0}$ à
$n_{7}$. %
Ces nœuds, reliés entre eux par des $0$-liaisons (en noir) et des $1$-liaisons
représentent une face. %
Puisque le membre gauche de la règle est vide, chaque nœud du membre droit est
créé. %
Puisque tous ces nœuds ont une étiquette vide $\langle \rangle$, chacun
d'entre eux crée un brin dans la G-carte, et par extension, une face «ex
nihilo». %
Les boucles de $2$-liaisons (bleu) et $3$-liaisons (verts) représentent le fait
que ces nœuds ne sont connectés à aucun autre nœuds par les liaisons $2$ et
$3$. %

\begin{figure}[h!]
  \centering
  \includegraphics[width=\textwidth]{rule-face-extrusion-std}
  \caption{Règle d'extrusion de face}
  \label{fig:rule-extrude-face}
\end{figure}

La règle d'extrusion de face (Figure~\ref{fig:rule-extrude-face}), elle,
comprend un unique nœud $n_{0}$ étiqueté $\langle 0,1 \rangle$ dans son membre
gauche et une $2$-liaison qui boucle sur $n_{0}$. %
Par le filtrage $m : n_{0} \mapsto a_{0}$ de $L$ vers $G$
(Figure~\ref{fig:face-to-extrude}), la face à extruder est entièrement
filtrée. %
Dans $H$ (Figure~\ref{fig:extruded-cube}), pour le membre droit, $n_{0}$
préserve la face extrudée (orange), $n_{1}$ crée des brins (bleu) copiés à
partir de ceux de $n_{0}$ (du membre gauche), $n_{2}$ crée des brins (magenta)
également copiés à partir de ceux de $n_{0}$ et ainsi de suite jusqu'au nœud
$n_{5}$. %
Les brins crées par les nœuds d'une règle, lorsque celle-ci à un membre gauche
non-vide, sont créés par copie des brins du nœud d'accroche. %
Ici, $n_{0}$ est le nœud d'accroche car c'est à partir de celui-ci que $m$
filtre un sous-graphe dans une G-carte. %

\begin{figure}[h!]
  \centering
  \includegraphics[width=.7\columnwidth]{rule-face-triangulation-3D-boxed-rapport.pdf}
  \caption{Règle de triangulation d'une face}
  \label{fig:rule-triangulation}
\end{figure}

La règle de triangulation de face (Figure~\ref{fig:rule-triangulation}) comporte
dans son membre gauche un unique nœud $n_{0}$ qui filtre une face, et dans son
membre droit, les nœuds $n_{0}$, $n_{1}$ et $n_{2}$. %
Cette règle lorsqu'elle est appliquée sur la face avant gauche de $H$, scinde
cette face est produit le graphe $I$
(Figure~\ref{fig:demo-cube-triangulated}). %
Dans le membre droit, $n_{0}$ a pour étiquette $\langle 0,\_ \rangle$, cela
signifie que les $1$-liaisons de la face sont supprimées lors de l'application
de la règle. %
Le nœud $n_{1}$, qui est donc créé par copie de $n_{0}$ est étiqueté
$\langle \_,2 \rangle$, ce qui signifie que dans cette copie, les $0$-liaisons
sont supprimées et les $1$-liaisons sont ré-étiquetées en $2$-liaisons. %
De plus, les brins de $n_{0}$ sont reliés à ceux de $n_{1}$ par des $1$-liaisons
et les brins de $n_{1}$ sont reliés à ceux de $n_{2}$ par des $0$-liaisons. %

Nous définissons deux types de liaisons dans une règle, d'une part les liaisons
dites explicites qui sont les liaisons entre les nœuds et d'une autre part les
liaisons dites implicites qui sont les liaisons représentées par les étiquettes
de nœuds. %
Notons que l'ordre des liaisons dans l'étiquette d'un nœud n'est pas anodins,
changer cet ordre revient à changer les liaisons entre les brins du sous-graphe
filtré. %

Tout comme dans les G-cartes, nous pouvons représenter des orbites directement
depuis les règles. %
Par exemple, l'orbite arête $\langle 0, 1, 3 \rangle$ incidente au nœud $n_{0}$,
dans le membre droit (encadrée en vert), comprend les nœuds $n_{0}$, $n_{1}$ et
$n_{2}$ car ceux-ci sont atteignables par leurs $0$-liaisons et $1$-liaisons
explicites. %
L'orbite sommet $\langle 1, 2, 3 \rangle$ incidente au nœud $n_{2}$, dans le
membre droit (encadrée en bleu), comprend uniquement le nœud $n_{2}$ car il n'y
a pas de liaison explicite étiquetée $1$, $2$ ou $3$ sortant de ce nœud. %

Les règles présentées ici, qui sont composées uniquement de nœuds, sont dites
«atomiques».%
En effet, Jerboa permet également de construire des opérations plus complexes
définies par des scripts, c'est-à-dire des enchaînements de règles. %

Enfin, les règles Jerboa sont munies de propriétés syntaxiques garantissant la
préservation de la cohérence des G-cartes, ainsi une G-carte de dimension 3 sera
toujours de dimension 3 après sa transformation par une règle. %

% \begin{figure}[h!]
%   \centering
%   \includegraphics[width=.8\columnwidth]{rule-sewA3}
%   \caption{Règle de coutures de deux faces par les $3$-liaisons}
%   \label{fig:rule-sew-A3}
% \end{figure}

% Par exemple, la règle (Figure~\ref{fig:rule-sew-A3}) qui assemble deux volumes
% prend en paramètre les deux faces le long desquelles les volumes seront
% cousus. %
% Ce type de règle, dite atomique, est composée de nœuds étiquetées décrivant des orbites dans une G-carte. %
% Les étiquettes de nœud sont appelées liaisons implicites, car elles représentes
% l'ensemble des liaisons de l'orbite décrite, tandis que les liaisons entre les
% nœuds sont dites liaisons explicites car elles décrivent les liaisons entre tous
% les brins de deux orbites. %


%%% Local Variables:
%%% mode: latex
%%% TeX-master: "../main"
%%% End:
